\section*{Limitations}
\label{sec:limitations}
\label{page:limitations_start}

\subsection*{Scale, stated as a falsifiable prediction}

The study operates at 5.6M parameters, $d_{\text{model}} = 256$, one block, a
vocabulary of 8{,}192, and three epochs of WikiText-103. Sparse-MoE benefits are
known to emerge with scale, so a reviewer can reasonably say the negative result
may be an artefact of operating below the regime where sparsity pays. \textbf{That
objection is correct and cannot be answered with the runs that exist.}

Rather than hedge it, we convert it into the paper's forward claim. The account in
\S\ref{sec:dilution} is that routing divides the \emph{diversity} of data each
parameter block sees, not its update budget---which the arithmetic there shows is
matched exactly. That account predicts a specific crossover: as the corpus grows at
fixed parameters, each expert's router-selected slice comes to contain enough of
the distribution's shared structure that the re-learning penalty amortises, so
\textbf{the MoR$-$MoRE gap should narrow monotonically with training data and the
ordering should eventually invert in favour of the composition}. The out-of-domain
result (\S\ref{sec:ood}) is weak early evidence in this direction, since MoRE is
already the better arm on data it was not trained on.

The measurement that would falsify it: hold parameters, depth budget and optimiser
fixed; sweep training tokens over at least a decade; and report the MoR$-$MoRE gap
as a function of tokens per expert. If the gap is flat or widening, the diversity
account is wrong and the composition's deficit is structural. We did not run that
sweep, and a second model size---the single most valuable follow-up---is likewise
unrun.

\subsection*{Budget, not convergence}\label{sec:budget} The 3-epoch budget leaves every arm short of convergence (\S\ref{sec:protocol}), and the arms are furthest apart on that axis at the cutoff: the final-epoch drop is $0.091$ nats for MoE, $0.132$ for MoR and $0.176$ for MoRE, so \textbf{the arm that loses in domain is the one still improving fastest}. The ordering is budget-dependent in the direction this paper's mechanism predicts---a routing arm that re-learns shared structure more slowly (\S\ref{sec:dilution}) should be behind at a fixed budget and closing---and we do not extrapolate a crossing epoch, because the curves decelerate. Appendix~\ref{app:budget} gives the per-epoch curves and the rest of this section.

\subsection*{Statistical resolution}

$n = 5$ seeds puts every headline comparison exactly on the design's resolution
floor of $p = 0.00794$ (\S\ref{sec:stats}). ``As extreme as a five-seed design
admits'' is the strongest statement available, and it is weaker than significance.
No larger effect could produce a smaller $p$ here; only more seeds could. Cohen's
$d$ values are reported alongside every gap precisely because a large $d$ on a
small gap is the misreading to avoid---across our own re-exports, $d$ has moved by
a factor of 1.5 while the underlying gap moved by less than 0.005 nats.

\subsection*{What was declared and not run}
\label{sec:declared}

One comparison in the pre-registered Family A, a depth-matched fixed-depth control,
was declared before results existed and \emph{still consumes its Holm slot}: we do
not delete it to loosen the thresholds on the survivors. The forced-depth sweep in
\S\ref{sec:halting} answers the same question using inference passes over the
canonical checkpoints rather than a trained control arm, so we report it as an
exploratory mechanism result and leave the declared slot consumed and unfilled. A
fixed-depth sweep trained from scratch, a shuffled-depth control, and a
depth--input mutual-information analysis were all declared exploratory and are
unrun.

\subsection*{Evaluation caveats}

BLiMP and EWoK are scored with the tokenizer fitted to our training corpus
(8{,}192 BPE merges), so \textbf{absolute accuracies are not comparable to
published models} using 32k--50k vocabularies; only the cross-arm comparison is
valid. EWoK sits within noise of its 50\% chance floor for every arm, which we
report as a scale limitation rather than a result---at 5.6M parameters a world
knowledge probe is below its useful range. All evaluations outside the two declared
families are uncorrected and carry no significance verdict in the abstract or
conclusions.

The perplexities in Table~\ref{tab:main} are intra-study quantities for the same
reason and should not be compared to published WikiText-103 numbers.

\subsection*{Disclosures}
\label{sec:honesty}

% T-LA.24. The label and the subsection heading STAY in the body: four places
% cite the disclosures in the Limitations section (02_setup protocol line, 02_setup GPU sentence,
% 03_method ablation pointer, 05_mechanism trapping aside), and a label that
% moves to the appendix would silently re-point all four. Only the prose moves.
% The replacement keeps every integrity-critical fact at body length: the
% specification violation, the partial pre-registration, the GPU/revision split
% with its measured common-mode offset, and the depth-metric N/A.
\textbf{Disclosures.} Three facts count against this study and we state them plainly. An earlier composed arm re-routed at every step, violated the specification, and its apparent parity with MoR was an artefact---the architecturally correct arm scores $0.134$ nats \emph{worse}. The pre-registration is partial, and the MoR-versus-MoE comparison was specified after its data existed. And the matrix spans two GPUs over four code revisions; the retrained arms reproduced to $-0.117\sigma$ and $-0.119\sigma$, a common-mode offset that cancels in every between-arm gap carrying a verdict (the MoE$-$MoR gap moved by 0.000185 nats), though per-seed values do not reproduce and we do not claim they do. We likewise report no per-token depth-allocation error: English supplies no ground-truth depth, so the metric is \texttt{N/A}, not zero. Appendix~\ref{app:disclosures} gives all of it in full.
